% Run GCT_KRON_20261120; GPT-6 Astra Pro; effort not exposed.
\section{A two-tree improvement of the joint later-row gate}

\subsection*{A-025. Simultaneously using the third and fifty-second rows}
In the BIP numerical range $n\geq m^{25}$, every actual
internally padded permanent multiplicity obstruction satisfies
\[
                   \boxed{\quad\lambda_3\geq557
                                  \quad\text{or}\quad\lambda_{52}\geq3.\quad}
                                                         \tag{A-025.1}
\]
This strengthens the sufficient constant $139$ in O-018.a.
The new construction was proposed by the root worker, and
the covering proof and exact uniform estimate were completed
here. No optimality among determinant representations is claimed.

\paragraph{A general joint-support cover.}
Assume $r\geq52$, and write $R=r-51$. Separate the
nonfirst variables into $x_2$, the forty-nine variables
$x_3,\ldots,x_{51}$, and the $R$ variables
$x_{52},\ldots,x_r$, called the early, middle, and last
groups in this proof. Suppose a coordinatewise downward-closed
support $\Gamma$ has total degree at most $m$, middle-plus-last
degree at most $B$, and last degree at most two. Other nested
conditions may be present. Let $u,v\geq0$ satisfy
\[
                              u+v=B-1,\qquad B\geq1.
\]
Define the following two monomial order ideals, intersected
with $\Gamma$:
\[
 \begin{split}
 P&=\Gamma\cap\{\text{early degree}\leq m,
                     \ \text{middle degree}\leq u,
                     \ \text{last degree}=0\},\\
 Q&=\Gamma\cap\{\text{early degree}=0,
                     \ \text{middle degree}\leq v,
                     \ \text{last degree}\leq1\}.
 \end{split}                                           \tag{A-025.2}
\]
Then $(P,Q)$ is admissible for F-018: with
$E=\{0,e_2,\ldots,e_r\}$ one has $\Gamma\subseteq P+Q+E$.
In particular
\[
 \kappa(\Gamma)\leq
       (m+1)\binom{49+u}{49}
          +(r-50)\binom{49+v}{49}-1.                  \tag{A-025.3}
\]

\begin{proof}
Take $\gamma\in\Gamma$. If its last degree is one or two,
remove one last-variable unit vector and put it on the
coefficient edge. Put the remaining last degree, at most one,
entirely in $Q$. The middle degree is then at most $B-1$.
If the last degree is zero but the middle degree is positive,
remove one middle-variable unit vector for the coefficient edge;
again the remaining middle degree is at most $B-1=u+v$.
Split this remaining middle multiset into degrees at most
$u$ and at most $v$, and put the two parts in $P,Q$,
respectively. Put every early-variable factor in $P$.
If both middle and last degrees are zero, put the whole early
monomial in $P$ and use $Q=0$ and edge exponent zero.

Each chosen endpoint is coordinatewise at most $\gamma$.
Downward closure therefore keeps both endpoints in $\Gamma$;
it also keeps every parent on their monomial-tree paths.
All requirements in (A-025.2) hold. This proves the cover.

Ignoring the intersections with $\Gamma$ can only enlarge
the cardinalities. Stars and bars gives
$|P|\leq(m+1)\binom{49+u}{49}$ and
$|Q|\leq(R+1)\binom{49+v}{49}$, since the last part
of $Q$ is either zero or one of its $R$ variables.
F-018 uses $|P|+|Q|-1$ vertices, proving (A-025.3).
The coefficient edge carries precisely the removed variable,
or $x_1$ when its exponent is zero. Thus the determinant
identity is the integral identity of F-018, with no change
to its cycle signs or padding exponents.
\end{proof}

\begin{proof}[A first uniform numerical bound]
The earlier actual-obstruction reductions used in O-018.a
give $m\geq190$ and $52\leq r\leq m^2$. Suppose
$\lambda_3\leq419$ and $\lambda_{52}\leq2$.
For the full nested support $\Gamma_{m,\lambda}$ of F-018,
the preceding cover applies with
\[
                         B=419,\qquad u=222,\qquad v=196.
\]
Write $C=\binom{271}{49}$ and $D=\binom{245}{49}$.
By (A-025.3), its determinant size is strictly less than
\[
                            (m+1)C+m^2D.
\]
The normalized upper bound is
\[
              \frac{(m+1)C+m^2D}{m^{25}}
                        =C(m^{-24}+m^{-25})+Dm^{-23}.
\]
Every term decreases as $m$ increases through positive real
values. At $m=190$, exact integer arithmetic gives
\[
       50\left(191\binom{271}{49}
                    +36100\binom{245}{49}\right)<49\cdot190^{25}.
                                                         \tag{A-025.4}
\]
The saved certificate expands both integers and their positive
difference; the binomials are computed by exact factorial
products, not floating point. Therefore for every $m\geq190$
the cover size is less than $(49/50)m^{25}<n$.

F-018 now gives $q^m_\lambda(d[n])\leq K_n(\lambda)$,
contradicting an actual obstruction. More explicitly, F-016
puts every coefficient generator of a determinant-ideal
highest-weight vector inside the nested support after aligning
the padding form. The cover represents every retained padded
form as a determinant, so every such vector vanishes on the
permanent closure. This includes every ideal multiplicity copy.
Thus the contradiction concerns the actual quotient dimensions.
It proves the weaker bound $\lambda_3\geq420$ or
$\lambda_{52}\geq3$. Retaining the total-degree constraint
below improves this to (A-025.1).
\end{proof}

\paragraph{Exact counts after total-degree pruning.}
The intersections in (A-025.2) retain total degree at most $m$.
Let $t=\min(m,u)$ and define
\[
 \begin{split}
 A(m,u)&=\sum_{j=0}^t\binom{48+j}{j}(m-j+1)\\
       &=(m+1)\binom{49+t}{49}-49\binom{49+t}{50},\\
 Q(m,v)&=\binom{49+\min(v,m)}{49}
        +(m^2-51)\binom{49+\min(v,m-1)}{49}.
 \end{split}                                             \tag{A-025.5}
\]
For $m\geq190$ and $r\leq m^2$, these give
\[
               \kappa(\Gamma)\leq A(m,u)+Q(m,v)-1.        \tag{A-025.6}
\]
Indeed a middle monomial of degree $j$ in $P$ has
$m-j+1$ possible early exponents, and there are
$\binom{48+j}{j}$ such middle monomials. The binomial sum
uses the identities
\[
 \sum_{j=0}^t\binom{48+j}{48}=\binom{49+t}{49},\qquad
 \sum_{j=0}^t j\binom{48+j}{48}=49\binom{49+t}{50},
\]
which follow from Pascal's identity and
$j\binom{48+j}{48}=49\binom{48+j}{49}$.
For $Q$, last degree zero allows middle degree through
$\min(v,m)$; last degree one has $R\leq m^2-51$
possible variables and allows middle degree through
$\min(v,m-1)$. This proves (A-025.5)--(A-025.6).
Other nested restrictions can only decrease these counts.

\begin{proof}[A first total-degree-pruned bound]
Suppose $\lambda_3\leq538$ and $\lambda_{52}\leq2$.
Use the same covering proof with
\[
                       B=538,\qquad u=325,\qquad v=212.
\]
These satisfy $u+v=B-1$. Set $T(m)=A(m,325)+Q(m,212)$.
The saved exact certificate verifies all 142 integer cases
\[
                  100T(m)<99m^{25}
                               \qquad(190\leq m\leq331).       \tag{A-025.7}
\]
Each record contains $m$, $T(m)$, $m^{25}$, and the positive
integer $99m^{25}-100T(m)$. The maximum ratio in this
finite range occurs at $m=331$, also checked by exact
cross multiplication. Every acceptance test in the certificate
uses integer arithmetic.

For the infinite tail $m\geq331$, write
$C'=\binom{374}{49}$ and $D'=\binom{261}{49}$.
The caps have then stabilized in (A-025.5), giving
\[
              T(m)=C'(m-635/2)+D'(m^2-50).
\]
For every real $m\geq331$, differentiation yields
\[
 2m^{26}\frac{d}{dm}\left(\frac{T(m)}{m^{25}}\right)
       =C'(15875-48m)+D'(2500-46m^2)<0.                 \tag{A-025.8}
\]
Both coefficients on the right are strictly negative there.
Thus the normalized count decreases on the whole tail,
and (A-025.7) at $331$ proves $T(m)<(99/100)m^{25}$
for every integer $m\geq331$ as well.

Combining the finite certificate and tail estimate shows
that the admissible cover has size less than $m^{25}\leq n$
for every allowed $m$. The same F-018 ideal inclusion used
in the first bound gives $q^m_\lambda\leq K_n(\lambda)$.
This contradiction proves $\lambda_3\geq539$ or
$\lambda_{52}\geq3$.
\end{proof}

\paragraph{Allocating the early variable to the other tree.}
There is a better cover in the same family. Put all early
degree in $Q$ instead of $P$. Thus $P$ has only middle
variables of degree at most $u$, while $Q$ has unrestricted
early degree, middle degree at most $v$, and last degree
at most one; both sets are intersected with $\Gamma$.
The covering proof is identical, with the early monomial
assigned to $Q$. Its total-degree-pruned bound is
\[
 \begin{split}
 \kappa(\Gamma)&\leq T_*(m;u,v)-1,\\
 T_*(m;u,v)&=\binom{49+\min(u,m)}{49}
                 +A(m,v)+(m^2-51)A(m-1,v),
 \end{split}                                            \tag{A-025.9}
\]
where $A$ is defined in (A-025.5). Last degree zero gives
the middle-plus-early count $A(m,v)$; last degree one
uses one of at most $m^2-51$ last variables and leaves
total degree $m-1$ for the other two groups.

\begin{proof}[Proof of (A-025.1)]
Suppose $\lambda_3\leq556$ and $\lambda_{52}\leq2$.
Use (A-025.9) with $u=360,v=195$, so $u+v=556-1$.
The certificate verifies all 171 integer cases
\[
             25T_*(m;360,195)<24m^{25}
                              \qquad(190\leq m\leq360).        \tag{A-025.10}
\]
The finite maximum ratio occurs at $m=360$; every
comparison is made by exact integer cross multiplication.

For real $m\geq360$, write $C''=\binom{409}{49}$,
$D''=\binom{244}{49}$, and $c=1911/10=49\cdot195/50$.
The normalized count is the sum of
\[
 \frac{C''}{m^{25}},\qquad
 \frac{D''(m+1-c)}{m^{25}},\qquad
 \frac{D''(m^2-51)(m-c)}{m^{25}}.                       \tag{A-025.11}
\]
Each is positive and strictly decreasing in this range.
For the last term its logarithmic derivative is
\[
 \frac{2m}{m^2-51}+\frac1{m-c}-\frac{25}{m}
       \leq\frac3m+\frac3m-\frac{25}{m}<0.
\]
Here $m^2\geq153$ implies the first bound, and
$m\geq360>3c/2$ implies $m-c\geq m/3$.
For the middle term the same bound gives logarithmic
derivative $1/(m+1-c)-25/m\leq3/m-25/m<0$.
The first term is plainly decreasing. Thus (A-025.10)
at $360$ proves the required bound for all larger integers.

The cover size is consequently less than $m^{25}\leq n$
for every $m\geq190$. F-018 again preserves all highest-weight
ideal copies and gives $q^m_\lambda\leq K_n(\lambda)$,
contradicting an actual obstruction. This proves (A-025.1).
\end{proof}

\paragraph{Verification and scope.}
The accompanying exact-integer script
\texttt{check\_joint\_support.py} and certificate
\texttt{JOINT\_SUPPORT\_CERTIFICATE.json} verify (A-025.4)
and every exact integer case in (A-025.7) and (A-025.10).
They also check all splits for the coarse bound at $m=190$:
$B=419$ has successful splits, while $B=420$ has none.
This last finite calculation is only about (A-025.3);
it is not an optimality assertion for the pruned support,
the full F-018 optimization, or arbitrary determinants.

The support uses the torus grading and exterior-minor argument
of F-016, with SGA 3, I, Proposition 4.7.3, pp.~35--36,
and Schur/exterior functors from SGA 6, VI, Theorem 3.3,
pp.~392--393. The determinant morphism and coefficient
identities are defined over $\mathbb Z$; the actual
highest-weight ideal comparison is in characteristic zero.
The argument uses two row constraints simultaneously and
retains the evaluation maps. It does not deduce the comparison
from a virtual lambda-ring identity or a support-saturation claim.
